\section{Broader impact}
\label{sec:impact}

Aggressive weight quantization lowers the compute and memory barriers to
running capable language models, which is a broadly positive development; it
also means that quality claims attached to low-bit releases can outrun what
users can verify. Independent forensics is one part of keeping that ecosystem
accountable: the ability to check, from public files, what a release actually
contains, how much of its quality is in its weights, and whether a baseline
survives off-calibration. Our findings are adversarial only in the scientific
sense---we redistribute no weights and rely on public artifacts---and we see no
direct negative societal impact specific to this study. The main risk we can
name is the inverse of its value: code-agreement or calibration-passage metrics
taken as quality guarantees. This paper documents two concrete ways that goes
wrong.

\section{Reproducibility and artifacts}
\label{sec:repro}

Every measurement in this paper is produced by scripts in the artifact, which
is available as \coderef{}. The artifact contains: the spec for the recovered
format with a canonical hash that is checked by the test suite
(\texttt{SPEC\_SHA256}, \code{0d2c008b4a9e\ldots}); byte-exact decoders for the
packed format; the rotation, layout, and agreement tooling; the controlled-swap
patcher; the Hessian capture and GPTQ variants; the evaluation harness wrapper;
offline tests that use synthetic tensors only (no downloads, CPU-only, CI
green); and the full failure register. Large artifacts---model weights,
calibration Hessians, evaluation outputs---are not redistributed; they are
regenerable from the public base and release with the commands in the artifact.
The evaluation protocol, hyperparameters, and the retraction timeline are
documented in Appendix~\ref{app:protocols}. We also note what is \emph{not}
reproducible: some provenance commit hashes refer to a private working history
and are provided for the record only.

\section{Conclusion}
\label{sec:conclusion}

We took apart a released ternary language model using only public files. The
format and the rotation basis are exactly recoverable from the release and the
public base; the last $8\%$ of codes are not, and they are the model: replacing
them with rotate-and-round weights of the base collapses perplexity by
$1{,}270\times$ in the same binary. Error-compensated quantization of the
public base cannot produce the released codes---every variant we tried moved
codes away, and the released codes fit real activation Hessians worse than
plain round-to-nearest---so the residual is weight movement from training, with
rotation-in-the-loop QAT/KD as the best-supported account and the mirror-map
recipe falsified in the deployed metric. On the evaluation side, a widely cited
public PTQ baseline reproduces its published $2.1\times$ only when calibration
and evaluation share a passage; off-calibration it collapses by three to four
orders of magnitude. We close the quantizer reverse-engineering track, retract
our own earlier 1.7B retention number after a holdout leak, and record every
falsified route so the next attempt starts from what was actually established:
the gap is localized, priced, and reachable only through end-to-end training at
scale.
